\subsection{Background}

A substation's protection and control equipment runs on DC power, usually from a
station battery and its charger. In normal service, the battery carries a steady
background load. When a fault occurs, protective relays energize the trip coils of the
breakers around the faulted equipment. Other operations may follow: breakers reclose,
a breaker-failure scheme trips neighboring breakers, and motor-operated mechanisms
recharge their springs. Each operation draws a large current for a fraction of a
second. The DC system must be able to supply the largest combination of these
\emph{momentary} currents that can occur at the same time~\cite{ieee946}.

Which equipment operates for a given fault depends on the \emph{protection zones}. A
zone is the part of the one-line diagram covered by one protection scheme, bounded by
the current transformers (CTs) that measure the current entering and leaving it. A
fault inside the zone trips the breakers that isolate it. Zones overlap at each
breaker so that no part of the station is left unprotected.

Finding the momentary load therefore means identifying every zone, the breakers and
motors that operate for a fault in it, the order in which they operate, and the DC
supply that feeds each coil. Done by hand, this involves tracing zones on paper
drawings and adding up coil currents in a spreadsheet. The work is slow, hard to
review, and must be repeated whenever the station changes.

\subsection{Project Objectives}

The goal of the project was a tool that lets the user:

\begin{itemize}
  \item transcribe a substation one-line diagram into an editable, connected drawing;
  \item edit the drawing's components and their electrical data;
  \item outline the zones of protection;
  \item calculate the momentary DC load of each operating stage; and
  \item identify the worst case: the largest momentary current on each DC supply.
\end{itemize}

The project also let the team apply what it learned in coursework, including power
system protection, software design, testing and troubleshooting, to a real
engineering problem for an industry sponsor.

\subsection{Scope}

\app{} meets these objectives. With it, the user can:

\begin{enumerate}
  \item show the station's one-line PDF as a background and rebuild the diagram from a
        library of electrical symbols;
  \item record each device's coil and motor currents, switch status, mechanism
        arrangement and DC supply;
  \item check the drawing for unconnected terminals and wiring shorts;
  \item have protection zones proposed automatically, then review, edit and verify them;
  \item calculate the momentary current for every zone, stage and DC supply, and the
        peak each supply must deliver; and
  \item save the study and export a PDF calculation report.
\end{enumerate}

Two areas are outside the tool's scope:

\begin{itemize}
  \item \textbf{Battery capacity.} The tool reports momentary current in amperes. It
        does not size the battery in ampere-hours, which requires a duty cycle over
        time (see IEEE Std~485~\cite{ieee485}).
  \item \textbf{Protection design.} The tool does not work out relay settings, CT
        assignments or trip logic. Its zones are proposals that an engineer confirms.
\end{itemize}

\subsection{Intended Users}

The tool is intended for GFT's substation protection and control engineers who prepare
or review DC load studies. Users should understand one-line diagrams and protection
zones. No programming knowledge is needed; the tool runs from a portable Windows
package.

\subsection{Report Organization}

\begin{itemize}
  \item Section~\ref{sec:overview} gives an overview of the tool: what it does, a
        typical workflow, the main window and how the software is built.
  \item Section~\ref{sec:installation} explains how to install and start it.
  \item Section~\ref{sec:quickstart} walks through a complete first study.
  \item Section~\ref{sec:userguide} is the user guide, organized by feature.
  \item Section~\ref{sec:shortcuts} lists every keyboard and mouse shortcut.
  \item Section~\ref{sec:calculation} explains the calculation, with a worked example.
  \item Section~\ref{sec:testing} describes how the tool was tested.
  \item Section~\ref{sec:future} covers limitations and future improvements, and
        Section~\ref{sec:conclusion} concludes.
  \item Appendix~\ref{app:formats} describes the file formats.
\end{itemize}
